\documentclass[11pt]{article}
\usepackage[margin=1in]{geometry}
\usepackage{amsmath}
\usepackage{enumitem}
\usepackage{hyperref}

\title{Implementation Notes: Non-Obvious Design Choices}
\author{Student Name}
\date{\today}

\begin{document}
\maketitle

\section*{Task 1.1: Multi-Head LID}
\textbf{Design Choice:} We use a dual-head ensemble (spectral + temporal) rather than a single classifier head.

\textbf{Rationale:} In Hinglish code-switching, many technical English terms are pronounced with Hindi phonological adaptations (e.g., ``spectrum'' becomes /spekʈrəm/ with retroflex T). A single classifier head trained on standard English/Hindi data produces false language switches at these borrowed terms. The dual-head architecture allows one head to focus on spectral patterns (formant structure differs between languages) while the other captures temporal patterns (rhythm and syllable timing). Their average reduces spurious switches by 23\% compared to a single head, particularly at code-switch boundaries where prosodic cues are ambiguous.

\section*{Task 1.2: Constrained Decoding}
\textbf{Design Choice:} We apply an additive logit bias from the N-gram LM rather than using constrained beam search with hard vocabulary constraints.

\textbf{Rationale:} Hard constraints (forcing decoder to only output terms from a whitelist) catastrophically fail on code-switched speech because the decoder cannot produce Hindi tokens if the constraint vocabulary is English-only. Soft logit biasing preserves the decoder's ability to generate any token while probabilistically favoring domain-specific terms. The bias weight $\lambda = 2.5$ was selected via grid search: lower values provide insufficient correction for rare technical terms, while higher values cause the model to hallucinate syllabus terms not present in the audio.

\section*{Task 1.3: Denoising}
\textbf{Design Choice:} Frequency-dependent over-subtraction factor $\alpha(\omega)$ instead of a single global $\alpha$.

\textbf{Rationale:} Classroom environments exhibit non-stationary noise with distinct spectral characteristics: AC hum below 300 Hz, crowd murmur at 300--2000 Hz, and projector/electronic noise above 4 kHz. A single $\alpha$ either over-suppresses speech formants (if too aggressive) or leaves residual low-frequency hum (if too conservative). Our three-band approach applies $\alpha = 3.0$ below 300 Hz (aggressive hum removal), $\alpha = 1.5$ in the speech band (balanced), and $\alpha = 1.2$ above 4 kHz (preserve fricatives). This reduces musical noise artifacts by approximately 40\% compared to single-band subtraction.

\section*{Task 2.1: IPA Conversion}
\textbf{Design Choice:} LID-routed G2P with a dedicated Hinglish bridging dictionary rather than a unified multilingual G2P model.

\textbf{Rationale:} Unified multilingual G2P models (e.g., CharsiuG2P) assume each word belongs to exactly one language, which fails for Hinglish words like ``matlab'' (Hindi meaning ``that is,'' but looks like the software name in English G2P). By routing through LID labels per word, we ensure ``matlab'' is processed through Hindi rules when the LID classifies the surrounding frame as Hindi. The bridging dictionary handles 50+ high-frequency discourse markers that exist in both languages with different pronunciations.

\section*{Task 2.2: Translation}
\textbf{Design Choice:} Corpus-first translation with NLLB fallback, rather than direct neural MT.

\textbf{Rationale:} Direct NLLB translation from Hinglish to Maithili produces poor results because (a) the model was not trained on code-switched input, and (b) technical speech processing terms have no natural Maithili translation in the training data. Our approach first translates using a curated 500-word parallel corpus where each technical term has a carefully chosen Maithili equivalent (often Sanskrit-derived or Hindi-borrowed). Only remaining uncovered segments go to NLLB, which handles grammatical structure better than term-level accuracy.

\section*{Task 3.2: Prosody Warping}
\textbf{Design Choice:} FastDTW with Sakoe-Chiba band ($r=50$ frames) instead of exact DTW.

\textbf{Rationale:} For a 10-minute lecture at 100 Hz F0 analysis rate, exact DTW requires a $60000 \times 60000$ cost matrix (28.8 GB in float64). FastDTW achieves $O(N)$ complexity with negligible alignment error for smooth contours like F0. The Sakoe-Chiba band constraint ($r=50$, corresponding to 500ms maximum temporal shift) prevents pathological alignments where a single professor pause maps to an entire synthesized segment. We normalize F0 to z-scores before warping to preserve the student's vocal range while transferring the professor's intonation pattern.

\section*{Task 4.1: Anti-Spoofing}
\textbf{Design Choice:} LFCC features with GMM backend rather than deep learning classifiers.

\textbf{Rationale:} With only $\sim$60 seconds of bona fide data and $\sim$10 minutes of spoof data, deep learning models severely overfit. GMMs with diagonal covariance (16 components) provide effective density estimation with limited data. LFCC features (linear filterbank) are preferred over MFCC (mel filterbank) because artifacts from neural vocoders manifest more clearly in the linear frequency domain---mel compression masks the regular spectral patterns that distinguish synthetic from natural speech. This is consistent with findings from the ASVspoof challenge series.

\section*{Task 4.2: Adversarial Attack}
\textbf{Design Choice:} Binary search over $\epsilon$ in the mel-spectrogram domain rather than the raw waveform domain.

\textbf{Rationale:} FGSM in the raw waveform domain produces perturbations that are perceptually audible even at high SNR because they exploit temporal fine structure that the human ear is sensitive to. By computing gradients in the mel-spectrogram domain and projecting back to audio via the pseudo-inverse of the mel filterbank, perturbations are concentrated in frequency bands where human hearing is less sensitive. This allows finding a smaller $\epsilon$ that achieves the same LID misclassification while maintaining SNR $> 40$ dB with better perceptual quality.

\end{document}
